\chapter{Disorientation}

\section*{Definition and Overview}
Disorientation is confusion about time, place, person, or situation. A person who is disoriented may not know what day it is, where they are, or who they or others are. Disorientation can be mild and brief, or severe and persistent, and it can result from many causes, including illness, infection, medicines, dehydration, head injury, dementia, and psychiatric conditions. Sudden disorientation is often a sign of a medical emergency such as delirium or stroke, and it warrants urgent assessment. This chapter explains the causes of disorientation, how it is assessed, and what to do.

\section*{Epidemiology}
Disorientation is common in hospitals, especially in older people, where it is a key feature of delirium, which affects up to a third of older hospitalised patients. It also occurs after surgery, with infections, and with medications. Dementia causes progressive disorientation, particularly to time and place in the later stages. Sudden disorientation in a previously alert person is a medical emergency and requires urgent evaluation.

\section*{Aetiology and Pathophysiology}
\begin{itemize}
    \item \textbf{Delirium:} acute confusion caused by infection, dehydration, medicines, surgery, pain, or other medical problems
    \item \textbf{Dementia:} progressive brain disease causing gradual disorientation, especially to time and place
    \item \textbf{Stroke:} sudden disorientation can result from a stroke affecting brain areas controlling awareness
    \item \textbf{Head injury:} concussion and brain injury cause confusion
    \item \textbf{Low blood sugar:} hypoglycaemia causes confusion and altered behaviour
    \item \textbf{Substances and medicines:} alcohol, drugs, and medication side effects
    \item \textbf{Metabolic and endocrine problems:} low sodium, liver or kidney failure, and thyroid problems
    \item \textbf{Psychiatric conditions:} severe depression, psychosis, and other mental illness
\end{itemize}

\section*{Clinical Features}
\begin{itemize}
    \item Confusion about the date, time, or season
    \item Not knowing where they are or how they got there
    \item Not recognising family members or not knowing their own name
    \item Difficulty following conversation or instructions
    \item Agitation, restlessness, or withdrawal
    \item Fluctuating alertness, especially in delirium
    \item Drowsiness, slow responses, or difficulty waking
\end{itemize}

\section*{Differential Diagnosis}
\begin{itemize}
    \item Delirium versus dementia, distinguished by onset and fluctuation
    \item Stroke and other neurological emergencies
    \item Metabolic causes such as hypoglycaemia, dehydration, and electrolyte imbalance
    \item Infection, especially urinary tract infection in older people
    \item Medication effects and substance intoxication or withdrawal
    \item Psychiatric conditions, including severe depression and psychosis
    \item Head injury, including in older people after falls
\end{itemize}

\section*{When to Seek Medical Care}
Any sudden or new disorientation requires urgent medical assessment, because it often indicates a serious underlying problem. Even mild confusion in an older person with a urinary infection or after a fall warrants prompt review.

\textbf{Call 000 (or local emergency services) for:}
\begin{itemize}
    \item Sudden confusion with weakness, facial droop, or difficulty speaking (possible stroke)
    \item Confusion with fever, drowsiness, or difficulty waking
    \item A head injury followed by confusion or memory loss
    \item Confusion with very low blood sugar or collapse
    \item Severe agitation, seizures, or inability to keep the person safe
\end{itemize}

\section*{Diagnosis and Investigations}
\begin{itemize}
    \item \textbf{History:} onset, duration, and circumstances, from the person and family
    \item \textbf{Cognitive assessment:} orientation to time, place, and person; attention and memory testing
    \item \textbf{Physical examination:} including temperature, hydration, and neurological signs
    \item \textbf{Blood tests:} glucose, electrolytes, kidney and liver function, blood count, and infection markers
    \item \textbf{Urine tests:} to detect urinary infection, a common cause in older people
    \item \textbf{Imaging:} CT or MRI of the brain for stroke, bleeding, or other structural causes
    \item \textbf{ECG and other tests:} where heart or metabolic causes are suspected
    \item \textbf{Medication review:} identifying drugs that cause confusion
\end{itemize}

\section*{Management}
\begin{itemize}
    \item \textbf{Treat the cause:} the underlying illness, infection, metabolic problem, or injury
    \item \textbf{Reassurance and orientation:} calm communication, reminding the person of time and place
    \item \textbf{Supportive care:} fluids, nutrition, and a familiar, quiet environment
    \item \textbf{Involve family:} familiar people help reorient and reassure
    \item \textbf{Reduce contributing medicines:} under medical supervision
    \item \textbf{Manage agitation:} non-drug approaches first; medication only when necessary
    \item \textbf{For dementia:} ongoing care and support from dementia services
\end{itemize}

\section*{Complications}
\begin{itemize}
    \item Wandering and getting lost, with risk of injury
    \item Falls and fractures
    \item Inability to care for oneself, with poor intake and dehydration
    \item Underlying emergencies such as stroke being missed
    \item Prolonged confusion and functional decline in older people
    \item Carer stress and need for increased supervision
\end{itemize}

\section*{Prognosis and Outcomes}
The outcome depends on the cause. Disorientation from reversible causes such as infection, dehydration, or low blood sugar resolves with treatment, usually within days. Delirium is reversible in most cases, although recovery can take weeks and some people do not return to their previous function. Disorientation from dementia is progressive, but good support improves safety and quality of life. Prompt assessment of sudden disorientation is essential for the best outcome.

\section*{Prevention and Self-Care}
\begin{itemize}
    \item Keep people with dementia oriented with calendars, clocks, and familiar routines
    \item Ensure glasses and hearing aids are used
    \item Prevent dehydration, poor nutrition, and infection in older people
    \item Supervise medications and review them regularly
    \item Prevent falls and head injuries
    \item Manage diabetes, blood pressure, and other conditions carefully
    \item Seek urgent care for sudden confusion
\end{itemize}

\section*{Sources and Further Reading}
The following reputable sources provide further information for healthcare students and patients seeking reliable, current advice about disorientation.

\begin{itemize}
    \item Healthdirect Australia. \textit{Confusion and disorientation: when to seek help.} https://www.healthdirect.gov.au/
    \item Dementia Australia. \textit{Disorientation and dementia support.} https://www.dementia.org.au/
    \item Australian Commission on Safety and Quality in Health Care. \textit{Delirium clinical care standard.} https://www.safetyandquality.gov.au/
    \item NHS. \textit{Confusion and delirium: symptoms and treatment.} https://www.nhs.uk/
    \item Mayo Clinic. \textit{Delirium: overview.} https://www.mayoclinic.org/
    \item MSD Manual. \textit{Confusion: professional overview.} https://www.msdmanuals.com/
\end{itemize}
